\documentclass[11pt,letterpaper,DIV=12]{scrartcl}

\usepackage[T1]{fontenc}
\usepackage[spanish,es-noshorthands]{babel}
\usepackage{lmodern}
\usepackage{graphicx}
\usepackage{booktabs}
\usepackage{tabularx}
\usepackage[hidelinks]{hyperref}
\usepackage{csquotes}
\usepackage[backend=biber,style=apa]{biblatex}

\addbibresource{references.bib}


\setcounter{secnumdepth}{0}

\begin{document}

\titlehead{\centering Instituto Tecnológico y de Estudios Superiores de Monterrey}
\subject{Inteligencia artificial avanzada para la ciencia de datos}
\title{Criterios de éxito de minería de datos}
\author{
  \begin{tabular}{l@{\hspace{2em}}l}
    Facundo Bautista Barbera & Emiliano Camacho Ponce \\
    A01066843 & A01712408 \\[1ex]
    Oswaldo Isaias Hernandez Santes & Iker Mejia Hernandez \\
    A01199004 & A01352358 \\[1ex]
    Alfredo Alejandro Soto Herrera & Jorge Manuel Oyoqui Aguilera \\
    A01711368 & A01711783
  \end{tabular}
}
\date{\today}

\maketitle


\section{Criterios objetivos}

\begin{itemize}
  \item \textbf{DM-1. Pronóstico de saturación.} El modelo tendrá éxito si detecta al menos el 80\% de los eventos en que algún robot supera el 90\% de utilización dentro de la ventana de 2 horas (recall $\geq$ 80\%), con una precisión de al menos 70\% en sus alertas. Esto equivale a aproximadamente una falsa alarma por robot por día como máximo, un nivel que el personal puede tolerar sin dejar de usar el sistema. Además, debe superar en F1 a un modelo base sencillo, como el promedio histórico de utilización por hora; de lo contrario, la complejidad del modelo no se justifica.
  \item \textbf{DM-2. Estimación de la congestión de la cola.} Como el VMS no registra la fila, el criterio es que la espera estimada a partir del tiempo ocioso entre ordeños coincida con la observación directa, con un error dentro de $\pm$20\% (o $\pm$5 minutos) en una muestra de horas pico. El segundo criterio es obtener un perfil horario de llegadas y de espera estimada para los 3 robots que caracterice el flujo circadiano y que el responsable del establo valide como representativo de la operación real.
  \item \textbf{DM-3. Estimación de la condición corporal.} El modelo tendrá éxito si su error absoluto medio (MAE) contra la calificación de los expertos del CAETEC es de 0.25 puntos o menos, con al menos 70\% de las predicciones dentro de $\pm$0.25 y al menos 90\% dentro de $\pm$0.50. Estas métricas deben sostenerse en un conjunto de prueba con vacas que no aparecieron en el entrenamiento, sin una degradación relevante respecto a validación. Para justificar la arquitectura Two-Stream, también debe mejorar el desempeño de un modelo equivalente que use solo RGB; si no lo hace, el canal LiDAR no está aportando valor.
  \item \textbf{DM-4. Detección de vacas en riesgo por BCS.} El sistema debe identificar al menos el 85\% de las vacas con BCS $\geq$ 4.00 al entrar al periodo seco o con BCS $\leq$ 2.25 (recall $\geq$ 85\%), porque omitir un caso tiene mayor costo que una falsa alarma, con una precisión de las alertas de al menos 70\%.
  \item \textbf{DM-5. Detección de monopolización y desplazamiento.} El análisis debe identificar al menos el 80\% de las vacas que, según la revisión del establo, quedan por debajo de la frecuencia de ordeño definida en horas pico (recall $\geq$ 80\%), de modo que el establo pueda actuar sobre las primíparas y las vacas de menor jerarquía.
  \item \textbf{DM-6. Detección de flujo bimodal.} El clasificador de sesiones bimodales debe alcanzar un F1 y un recall de al menos 0.80 frente a la revisión manual de una muestra de perfiles de flujo, para que la estimación de leche recuperable sea confiable.
\end{itemize}

\section{Criterios subjetivos}

Además de las métricas, el éxito depende de tres juicios. El responsable del establo debe considerar que las alertas de saturación llegan con anticipación suficiente y son claras para decidir cómo rebalancear el tráfico. El médico veterinario debe considerar que la calificación automática de BCS es tan útil como la evaluación visual para decidir ajustes de nutrición y de manejo reproductivo. El socio formador debe confirmar que la solución puede operarse con el personal y el equipo actuales (iPhone y registros del VMS) sin interrumpir la rutina de ordeño.

\section{Criterio mínimo de éxito}

Dado que los indicadores productivos y de salud requieren semanas de operación para mostrar cambios, al cierre en diciembre de 2026 el éxito mínimo consiste en tres cosas. Primero, que DM-3 cumpla el MAE y los porcentajes de acierto definidos en el conjunto de prueba y que DM-4 detecte los casos de riesgo con el recall establecido. Segundo, que con datos históricos del CAETEC DM-1 demuestre que habría anticipado los periodos de saturación registrados, con el recall definido, y que se estime su efecto sobre los indicadores de negocio. Tercero, que se realice al menos una prueba de captura con iPhone en el establo, validada por el médico veterinario y el responsable del establo. La verificación completa de las metas productivas y de salud se propone para un piloto posterior de al menos 8 semanas, acordado con el CAETEC.

\printbibliography

\end{document}
